\documentclass[11pt]{article}
\usepackage[T1]{fontenc}
\usepackage{lmodern,amsmath,amssymb,amsthm,mathtools,booktabs,microtype}
\usepackage[margin=1in]{geometry}
\usepackage{xurl}
\usepackage[colorlinks=true,linkcolor=blue,urlcolor=blue,citecolor=blue]{hyperref}
\hypersetup{pdftitle={A residue-field optimal test for symplectic blocks at every congruence precision},pdfauthor={Research report prepared for Vladimir Reshetnikov}}
\newtheorem{theorem}{Theorem}
\newtheorem{lemma}[theorem]{Lemma}
\newtheorem{proposition}[theorem]{Proposition}
\newtheorem{corollary}[theorem]{Corollary}
\newtheorem{remark}[theorem]{Remark}
\newcommand{\OO}{\mathcal O}
\newcommand{\F}{\mathbb F}
\newcommand{\Pp}{\mathbb P}
\newcommand{\E}{\mathbb E}
\newcommand{\Alt}{\operatorname{Alt}}
\newcommand{\Mat}{\operatorname{Mat}}
\newcommand{\Sp}{\operatorname{Sp}}
\newcommand{\TV}{d_{\mathrm{TV}}}
\newcommand{\Pf}{\operatorname{Pf}}
\newcommand{\rank}{\operatorname{rank}}
\newcommand{\GL}{\operatorname{GL}}
\newcommand{\length}{\operatorname{length}}
\newcommand{\qbin}[3]{\genfrac{[}{]}{0pt}{}{#1}{#2}_{#3}}
\title{A residue-field optimal test for symplectic blocks\newline at every congruence precision}
\author{Research report prepared for Vladimir Reshetnikov}
\date{27 November 2026}
\begin{document}
\maketitle
\begin{abstract}
Let a block be observed in a Haar-random symplectic matrix over the integers of a finite extension of $\mathbb Q_p$. Suppose its source and target are nondegenerate and the orthogonal complement of the target has at least the source dimension. We prove that the optimal equal-prior test against independent Haar entries uses only the first residue digit, for every source dimension and every precision. The event is singularity of one residual alternating form. We give a closed finite formula for the total variation distance, valid also in residue characteristic two. The proof combines the classical alternating local-density formula with a sharp sign inequality and an elementary Fourier calculation over the residue field. The result extends the one-pair theorem in the inspected ProveIt article and answers its multi-pair precision-persistence question in this dimension range. Higher digits can still improve tests with unequal priors; this is proved explicitly.
\end{abstract}

\section{Statement and relation to the inspected source}
The inspected source is \texttt{Analysis/GeneralizedRiemannHypothesis/Article/sections/spn.tex} in ProveIt, at commit\newline
\texttt{2e714ce3d8e33aa8160391cf5dd95f3af3c34045} \cite{repo}.

Its subsection \texttt{spn:sec:padic} proves precision-independent total variation when the source is one symplectic pair. Its final question \texttt{spn:q:symplectic} asks for the optimal local test when both frames contain several symplectic pairs. The theorem below treats every number of pairs, assuming a complement large enough to contain the source. It is an unconditional extension of a method in that directory; it uses no assertion about zeros of an $L$-function.

Let $K/\mathbb Q_p$ be a finite extension, $\OO$ its integer ring, $\pi$ a uniformizer, and $\OO/\pi\OO=\F_q$. Put $R_e=\OO/\pi^e\OO$ and $R_\infty=\OO$. We allow $p=2$ and ramified extensions. Let $J_j$ be the block diagonal matrix with $j$ blocks $\left(\begin{smallmatrix}0&1\\-1&0\end{smallmatrix}\right)$. Fix integers
\[
 a\ge1,\qquad r\ge1,\qquad c\ge r,\qquad n=a+c.
\]
A uniform symplectic $2r$-frame in $R_e^{2n}$ can be written
\[
 X=\begin{pmatrix} A\\T\end{pmatrix},\qquad
 A\in\Mat_{2a,2r}(R_e),\quad T\in\Mat_{2c,2r}(R_e),\qquad
 X^{\mathsf T}J_nX=J_r.
\]
Equivalently, $X$ comprises a fixed collection of symplectic column pairs of a Haar-random element of $\Sp_{2n}(R_e)$. Let $\mu_e$ be the law of the observed block $A$ and let $\nu_e$ be the law of $4ar$ independent uniform entries. Both finite and infinite Haar measures are normalized to mass one. Define
\[
 S(A)=J_r-A^{\mathsf T}J_aA,
 \qquad E=\{A:\Pf(S(A))\equiv0\pmod\pi\}.
\]
For $v\ge r$ and $j\ge0$ set
\[
 \beta_{v,r}=\prod_{i=v-r+1}^{v}(1-q^{-2i}),\qquad
 \eta_j=\prod_{i=1}^j(1-q^{-2i}),\qquad
 \theta_j=\prod_{i=1}^j(1-q^{-(2i-1)}),
\]
with empty products one. The Gaussian binomial with parameter $Q$ is
\[
 \qbin rsQ=\frac{\prod_{i=1}^r(1-Q^i)}{
  \prod_{i=1}^s(1-Q^i)\prod_{i=1}^{r-s}(1-Q^i)}.
\]

\begin{theorem}[All-pairs precision invariance]\label{thm:main}
Under these assumptions, for every $e\in\{1,2,\ldots,\infty\}$,
\begin{equation}\label{eq:tv}
 \boxed{\displaystyle
 \TV(\mu_e,\nu_e)=
 \left(1-\frac{\beta_{c,r}}{\beta_{n,r}}\right)
 \sum_{s=0}^{r}\qbin r{s}{q^{-2}}\theta_{r-s}q^{-(2a+1)s}.}
\end{equation}
The event $E$ is optimal for distinguishing the two distributions with equal priors: accept the symplectic-block model exactly on $E$. The likelihood is strictly below one off $E$ and strictly above one on $E$, up to null sets. Off $E$, it is the constant $\beta_{c,r}/\beta_{n,r}$. The minimum equal-prior error probability is $(1-\TV(\mu_e,\nu_e))/2$ and is therefore unchanged by revealing any further digits.
\end{theorem}

The same statement holds for arbitrary fixed unimodular source and target alternating forms, after changing symplectic bases. If one observes the pairing matrix $M=Y^{\mathsf T}J_n gX_0$ instead of the coordinate block, the residual form is
$H_0+M^{\mathsf T}H_1^{-1}M$. Fixed invertible changes of variables preserve the independent-entry Haar law.

For $r=1$, the sum in \eqref{eq:tv} is $1-q^{-1}+q^{-(2a+1)}$, and
\[
 1-\frac{\beta_{c,1}}{\beta_{n,1}}=\frac{q^{2a}-1}{q^{2n}-1}.
\]
Thus the theorem recovers the article's one-pair formula exactly. For $r=2$, the sum is
\[
 (1-q^{-1})(1-q^{-3})+(1+q^{-2})(1-q^{-1})q^{-(2a+1)}+q^{-2(2a+1)}.
\]
For example, $q=2,a=r=c=2$ gives $\nu(E^c)=469/1024$ and
$\TV(\mu_e,\nu_e)=6767/52224\approx0.1295764400$ at every precision.

\section{The classical density input and its normalization}
For independent Haar entries $T\in\Mat_{2v,2r}(\OO)$, write
$G_v=T^{\mathsf T}J_vT$. For a nondegenerate alternating matrix $C$ over $K$ with integral entries, its alternating elementary-divisor type is the partition
$\lambda=(\lambda_1\ge\cdots\ge\lambda_r\ge0)$ determined by
\[
 C\sim\operatorname{diag}(\pi^{\lambda_1}J_1,\ldots,\pi^{\lambda_r}J_1)
\]
under integral congruence. Let $B$ be any complete discrete valuation ring of residue cardinality $q^2$ and put
$M_\lambda=\bigoplus_{i=1}^r B/\pi^{\lambda_i}B$. The polynomial
\[
 P_\lambda(z;q^2)=\sum_{N\le M_\lambda}z^{\length_B N}
\]
counts all submodules of this finite module, including zero. It depends only on $q$ and $\lambda$, and has nonnegative integer coefficients. Its constant coefficient is one. If $\lambda\ne0$, it has a term of degree one, so
\begin{equation}\label{eq:poly-lower}
 P_\lambda(z;q^2)\ge1+z\qquad(z>0).
\end{equation}
The notation $B$ and its uniformizer in this definition is auxiliary: no quadratic extension of $K$ is part of the observation model.

\begin{lemma}[Classical Gram density]\label{lem:gram}
For $v\ge r$, $G_v$ has a density relative to Haar measure on $\Alt_{2r}(\OO)$. On the orbit of a nondegenerate $C$ of type $\lambda$, it equals
\begin{equation}\label{eq:gram}
 H_{v,r}(C)=\beta_{v,r}
 P_\lambda\bigl(q^{-(2v-2r+1)};q^2\bigr).
\end{equation}
In particular it equals $\beta_{v,r}$ when $C$ is unimodular.
\end{lemma}

This is a classical input, not a new density formula. A directly usable primary reference is Wakatsuki \cite[Theorem 4.3 and Proposition 4.4]{wakatsuki}. Here is the precise normalization check, which also avoids an ambiguity in the printed definition of the Hall polynomial on page 467 of that paper. Use his $n=v,m=r,l=2r$. Let $f^\lambda_{\mu\nu}(t)$ be the structure constants of the usual Hall--Littlewood polynomials and let $n(\lambda)=\sum_i(i-1)\lambda_i$. Dividing his Proposition 6.7 by the orbit volume in Lemma 4.1 gives
\[
 \frac{\eta_v}{\eta_{v-r}}
 \sum_{\mu,\nu}
 q^{2n(\lambda)-2n(\mu)-2n(\nu)}
 f^\lambda_{\mu\nu}(q^{-2})
 q^{-(2v-2r+1)|\mu|}.
\]
The coefficient
$q^{2n(\lambda)-2n(\mu)-2n(\nu)}f^\lambda_{\mu\nu}(q^{-2})$
is the standard Hall number $g^\lambda_{\mu\nu}(q^2)$. It counts submodules of type $\nu$ with quotient of type $\mu$ in a module of type $\lambda$ over residue field $q^2$. Duality of finite modules, or the symmetry $g^\lambda_{\mu\nu}=g^\lambda_{\nu\mu}$, shows that summing by quotient length gives the same polynomial as summing by submodule length. This gives \eqref{eq:gram}. The older general alternating-density formulas are due to Hironaka and Sato \cite{hironaka-sato}.

\begin{lemma}[A strict sign gap]\label{lem:gap}
For all $q\ge2$, $r\ge1$ and $v\ge r$,
\[
 \beta_{v,r}\left(1+q^{-(2v-2r+1)}\right)>1.
\]
Consequently $H_{v,r}(C)<1$ on unimodular $C$, and $H_{v,r}(C)>1$ on every other nondegenerate $C$.
\end{lemma}
\begin{proof}
Put $d=v-r$, $y=q^{-2d}$ and $x=y/q$. The finite product bound gives
\[
 \beta_{v,r}\ge1-\sum_{i=d+1}^{d+r}q^{-2i}
 >1-\frac{y}{q^2-1}.
\]
The strict inequality is valid even for $r=1$, because the omitted infinite tail is positive. Now
\[
 (1+x)\left(1-\frac{y}{q^2-1}\right)-1
 =\frac{y(q^2-q-1-y)}{q(q^2-1)}\ge0,
\]
since $q\ge2$ and $0<y\le1$. The strict product bound proves the claim. Combine it with \eqref{eq:poly-lower}; a nonunimodular nondegenerate $C$ has $\lambda\ne0$.
\end{proof}

\section{From the density to an optimal residue test}
For finite $e$, let $H_{v,r,e}(C)=q^{eD}\Pp(G_v\equiv C\bmod\pi^e)$, where $D=\binom{2r}{2}$. A frame of Gram matrix $J_r$ is automatically primitive and extends to a symplectic basis over $R_e$. The Gram map is smooth at every such frame: its differential
$Z\mapsto X^{\mathsf T}J_vZ+Z^{\mathsf T}J_vX$ is onto the alternating matrices, because $X$ has a left inverse respecting its unimodular Gram matrix. This assertion uses no division by two. Counting over the residue field and lifting therefore gives
\[
 \#\{X\in\Mat_{2v,2r}(R_e):X^{\mathsf T}J_vX=J_r\}
 =q^{e(4vr-D)}\beta_{v,r}.
\]
Splitting the first $2a$ coordinates from the complement now proves the exact finite likelihood
\begin{equation}\label{eq:finite-likelihood}
 h_e(A):=\frac{\mu_e(A)}{\nu_e(A)}
 =\frac{H_{c,r,e}(S(A))}{\beta_{n,r}}.
\end{equation}

On the $p$-adic matrix space the corresponding density is, almost everywhere,
\begin{equation}\label{eq:likelihood}
 h_\infty(A)=\frac{H_{c,r}(S(A))}{\beta_{n,r}}.
\end{equation}
For completeness, no singular mass is lost in passing to this formula. On the open set where $\Pf(S(A))\ne0$, the alternating elementary-divisor orbit is open, so for each $A$ the finite density in \eqref{eq:finite-likelihood} equals \eqref{eq:likelihood} throughout a sufficiently small neighborhood once the precision is large enough. The complement has $\nu_\infty$-measure zero, since $\Pf(S(A))$ is a nonzero polynomial (at $A=0$ it equals one). It also has $\mu_\infty$-measure zero: on the symplectic frame variety, it is a nonzero regular function, since a frame contained entirely in the complement has residual $J_r$. That variety is a smooth irreducible homogeneous space. In its rational coordinate charts, a nonzero polynomial has Haar-null zero set, by induction and Fubini; hence the same holds on the frame variety. The two laws therefore have density relation \eqref{eq:likelihood} everywhere except null sets.

Since $a\ge1$, one has $0<\beta_{c,r}<\beta_{n,r}<1$. On $E^c$ the likelihood is the constant
$\alpha=\beta_{c,r}/\beta_{n,r}<1$. On $E$, Lemmas \ref{lem:gram}--\ref{lem:gap} give $h_\infty>1$ almost everywhere. The finite density is its conditional expectation given $A\bmod\pi^e$. Every such residue class lies either in $E$ or in $E^c$, so the same strict sign holds at every precision. Thus
\begin{equation}\label{eq:tv-reduction}
 \TV(\mu_e,\nu_e)=\int_{E^c}(1-h_e)\,d\nu_e
 =(1-\alpha)\nu_1(E^c).
\end{equation}
This proves precision invariance; the remaining task is to evaluate $\nu_1(E^c)$ explicitly.

\section{An elementary finite-field calculation}
The nonsingular residual probability evaluated below is a known finite-field counting ingredient. It is also obtained from De Boeck and Van de Voorde \cite[Theorem 4.6]{deboeck}, specialized to $i=0,j=r,n=a+r$ and divided by $q^{4ar}$, after converting the Gaussian-binomial parameter from $q^2$ to $q^{-2}$. Thus the probability polynomial itself is not claimed as a new counting theorem. We include a short Fourier proof to keep its use and normalization explicit.
All calculations here hold in every finite field, including characteristic two. The pairing on alternating matrices is
$\langle W,C\rangle=\sum_{i<j}W_{ij}C_{ij}$; there is no factor $1/2$. Fix a nontrivial additive character $\psi$ of $\F_q$.

The Fourier eigenvalue in the next lemma is likewise classical, as the full-rank relation of the alternating-forms association scheme of Delsarte and Goethals \cite{delsarte}; a short direct proof is supplied.
\begin{lemma}[Fourier transform of the nonsingular orbit]\label{lem:fourier}
If $W\in\Alt_{2r}(\F_q)$ has rank $2s$, then
\begin{equation}\label{eq:fourier}
 F_r(s):=\sum_{\substack{C\in\Alt_{2r}(\F_q)\\C\text{ nonsingular}}}
 \psi(\langle W,C\rangle)
 =(-1)^s q^{r(r-1)}\prod_{i=1}^{r-s}(q^{2i-1}-1).
\end{equation}
\end{lemma}
\begin{proof}
Congruence invariance reduces $W$ to $s$ standard symplectic blocks followed by zeros. For $s=0$ the sum is the number of nonsingular alternating matrices,
\[
 F_r(0)=\frac{|\GL_{2r}(\F_q)|}{|\Sp_{2r}(\F_q)|}
 =q^{r(r-1)}\prod_{i=1}^{r}(q^{2i-1}-1).
\]
For $s\ge1$, write
\[
 C=\begin{pmatrix}0&t&u\\-t&0&v\\-u^{\mathsf T}&-v^{\mathsf T}&D\end{pmatrix}.
\]
If $D$ is singular, expansion of the Pfaffian shows that nonsingularity of $C$ is independent of $t$: the coefficient of $t$ is $\Pf(D)=0$. Summing $\psi(t)$ cancels this contribution. If $D$ is nonsingular, its Schur complement says that precisely one value of $t$ is forbidden. The sum over the other values is minus the character of that forbidden value. Summing next over $u,v$ gives $q^{2r-2}$, by orthogonality for the nonsingular pairing with matrix $D^{-1}$. The character involving $D$ has rank $2(s-1)$. Therefore
$F_r(s)=-q^{2r-2}F_{r-1}(s-1)$, and induction proves \eqref{eq:fourier}. The case $r=1$ uses the empty matrix and the same convention.
\end{proof}

Let $A_s$ denote the number of alternating $2r$ by $2r$ matrices of rank $2s$. Choosing their radical and a nondegenerate form on the quotient gives
\begin{equation}\label{eq:rank-count}
 A_s=q^{s(4r-2s-1)}
 \frac{\gamma_{2r}}{\gamma_{2r-2s}\eta_s},
 \qquad\gamma_j=\prod_{i=1}^j(1-q^{-i}).
\end{equation}
In particular $A_r=q^D\theta_r$.

\begin{proposition}[Nonsingular residual probability]\label{prop:prob}
If $A\in\Mat_{2a,2r}(\F_q)$ has independent uniform entries, then
\begin{equation}\label{eq:prob}
 \Pp(J_r-A^{\mathsf T}J_aA\text{ nonsingular})
 =\sum_{s=0}^r\qbin r{s}{q^{-2}}\theta_{r-s}q^{-(2a+1)s}.
\end{equation}
\end{proposition}
\begin{proof}
Group the $2a$ rows of $A$ into $a$ independent pairs. The characteristic function of each alternating outer-product difference at $W$ is $q^{-\rank W}$, by averaging one vector and then the other. Thus the Gram characteristic function is $q^{-a\rank W}$.

Let $K_s=\sum_{\rank W=2s}\psi(\langle W,J_r\rangle)$. Summing the same character over the pair of rank orbits in the two possible orders gives
$A_rK_s=A_sF_r(s)$. Fourier inversion, applied to the indicator of the nonsingular orbit, yields
\[
 \Pp(J_r-A^{\mathsf T}J_aA\text{ nonsingular})
 =q^{-D}\sum_{s=0}^r q^{-2as}\frac{A_s}{A_r}F_r(s)^2.
\]
The orbit reciprocity is valid for the dual congruence action as well as the usual congruence action; the contragredient sends rank orbits to rank orbits. There is no sign issue, because all rank orbits are stable under negation.

Now $F_r(s)=(-1)^s q^{D-s(2r-s)}\theta_{r-s}$ and $\gamma_{2j}=\eta_j\theta_j$. Substitution in \eqref{eq:rank-count} simplifies the coefficient of $q^{-2as}$ to
\[
 q^{-D}\frac{A_s}{A_r}F_r(s)^2
 =q^{-s}\frac{\eta_r}{\eta_s\eta_{r-s}}\theta_{r-s}
 =q^{-s}\qbin r{s}{q^{-2}}\theta_{r-s}.
\]
This is \eqref{eq:prob}. Together with \eqref{eq:tv-reduction}, it proves Theorem \ref{thm:main}.
\end{proof}

\section{What precision invariance does and does not imply}
Theorem \ref{thm:main} concerns one block, two specified distributions, equal priors and zero-one loss. The full likelihood is not determined by its first digit. Thus precision invariance is a statement about the sign of the likelihood minus one, not statistical sufficiency.

\begin{proposition}[Higher digits can improve unequal-prior decisions]\label{prop:priors}
For every parameter choice in Theorem \ref{thm:main}, there are unequal priors for which observing three digits gives a strictly smaller Bayes error than observing only the first digit.
\end{proposition}
\begin{proof}
Take $A(t)$ whose first two columns in the first target plane are $(e_1,t e_2)$ and all other entries are zero. Its residual is
$\operatorname{diag}((1-t)J_1,J_{r-1})$. For $t=1-\pi$ and $t=1-\pi^2$, both blocks have the same reduction modulo $\pi$, while their alternating elementary-divisor types are respectively $(1,0,\ldots,0)$ and $(2,0,\ldots,0)$. These types persist in open neighborhoods modulo $\pi^3$, each of positive $\nu_\infty$-measure.

Put $x=q^{-(2c-2r+1)}$ and $\alpha=\beta_{c,r}/\beta_{n,r}$. Since a cyclic module has one submodule of every possible length, the two likelihoods on these neighborhoods are
\[
 L_1=\alpha(1+x),\qquad L_2=\alpha(1+x+x^2),\qquad L_1<L_2.
\]
Choose $L_1<\tau<L_2$ and prior probability $\rho=1/(1+\tau)$ for the symplectic model. The optimal test has likelihood threshold $\tau$. The first-digit residue class just constructed contains sets of positive probability on both sides of the threshold, whereas three digits separate those two neighborhoods. Equivalently, the conditional triangle inequality for $|\rho h-(1-\rho)|$ is strict on that class. Hence Bayes error strictly decreases. Lemma \ref{lem:gap} shows $L_1>1$, so these priors are indeed unequal.
\end{proof}

The dimension condition also matters. If $c=0$, the residual form must be exactly zero and the $p$-adic block law is singular relative to independent Haar entries. The theorem does not assert persistence for $0<c<r$, for degenerate source or target forms, or for multiple independent observations. A product likelihood can cross one even when every individual sign is already visible at the first digit.

\section{Verification, contribution and limitations}
The supplied script performs independent finite-ring checks using character orthogonality, not the Hall-polynomial density formula. For $4$ by $4$ alternating matrices over $\mathbb Z/p^e\mathbb Z$, it enumerates every Fourier character and obtains exact rational probabilities by averaging over unit scalars (Ramanujan sums). It checks normalization, all alternating elementary-divisor classes, the density formula whenever the invariant factors are resolved, and precision-invariant Gram total variation. It also enumerates small observed blocks and checks the full compression likelihood, total variation formula and optimal event. The output file records the exact parameters and counts: 13 Gram parameter sets, 11 compression parameter sets, 1,863,252 Fourier characters and 407,106 observed blocks. All checks pass.

These finite checks are corroboration. The proof for all dimensions is the local-density theorem plus the explicit sign and Fourier arguments above. No Lean formalization of the new theorem is claimed. The primary density theorem is stated for integer rings of $p$-adic number fields, and the present theorem deliberately keeps that scope; no equal-characteristic extension is needed.

Relative to the inspected article, the new contribution is the all-pairs optimal first-digit test, its exact total variation formula, and the explicit separation between equal-prior testing and statistical sufficiency. The classical local density and standard Hall-number interpretation are credited as inputs. A focused literature search located those density results and did not identify this precise statistical consequence; this is not a claim to have certified worldwide priority.

\begin{thebibliography}{9}
\bibitem{delsarte} P. Delsarte and J. M. Goethals, \emph{Alternating bilinear forms over $\mathrm{GF}(q)$}, Journal of Combinatorial Theory, Series A \textbf{19} (1975), 26--50.\newline
\url{https://doi.org/10.1016/0097-3165(75)90090-4}.
\bibitem{deboeck} M. De Boeck and G. Van de Voorde, \emph{Anzahl theorems for trivially intersecting subspaces generating a non-singular subspace I: symplectic and hermitian forms}, Linear Algebra and its Applications \textbf{699} (2024), 367--402.\newline
\url{https://doi.org/10.1016/j.laa.2024.07.004}. Open primary text:\newline
\url{https://arxiv.org/html/2407.07486v3}.
\bibitem{repo} V. Reshetnikov, \emph{ProveIt: Generalized Riemann Hypothesis research article}, section \texttt{spn.tex}, commit \texttt{2e714ce3d8e33aa8160391cf5dd95f3af3c34045}, inspected 27 November 2026.\newline
\url{https://github.com/VladimirReshetnikov/ProveIt/blob/2e714ce3d8e33aa8160391cf5dd95f3af3c34045/Analysis/GeneralizedRiemannHypothesis/Article/sections/spn.tex}
\bibitem{wakatsuki} S. Wakatsuki, \emph{Igusa local zeta functions and integration formulas associated to equivariant maps}, Osaka Journal of Mathematics \textbf{42} (2005), 463--486. In particular Lemma 4.1, Theorem 4.3, Proposition 4.4 and Proposition 6.7.\newline
\url{https://doi.org/10.18910/4504}. Open primary text:\newline
\url{https://ir.library.osaka-u.ac.jp/repo/ouka/all/4504/1704ojm.pdf}
\bibitem{hironaka-sato} Y. Hironaka and F. Sato, \emph{Local densities of alternating forms}, Journal of Number Theory \textbf{33} (1989), 32--52.\newline
\url{https://doi.org/10.1016/0022-314X(89)90058-9}.
\end{thebibliography}
\end{document}
